% Springer-compatible LaTeX template for APJM / MIR submission
% Compiled from Markdown via pandoc.
% Document class: `article` (Springer Nature's sn-jnl is not bundled; this template
% follows Springer's general manuscript guidelines and is compatible with both
% Springer Nature LaTeX template and Word-track submission portals).

\documentclass[12pt,a4paper]{article}

% --- Geometry & spacing ---
\usepackage[margin=2.5cm]{geometry}
\usepackage{setspace}
\onehalfspacing

% --- Encoding & fonts ---
% Conditional setup: under pdflatex use legacy inputenc; under xelatex/lualatex
% use fontspec for native Unicode support (needed for β, −, × symbols).
\usepackage{iftex}
\ifPDFTeX
  \usepackage[utf8]{inputenc}
  \usepackage[T1]{fontenc}
  \usepackage{lmodern}
\else
  \usepackage{fontspec}
  \setmainfont{Liberation Serif}[
    Scale=1.0,
    BoldFont={Liberation Serif Bold},
    ItalicFont={Liberation Serif Italic},
    BoldItalicFont={Liberation Serif Bold Italic}
  ]
  \setmonofont{Liberation Mono}[Scale=0.85]
  % Map common Unicode math/typographic symbols
  \usepackage{newunicodechar}
  \newunicodechar{∈}{\ensuremath{\in}}
  \newunicodechar{∉}{\ensuremath{\notin}}
  \newunicodechar{≤}{\ensuremath{\leq}}
  \newunicodechar{≥}{\ensuremath{\geq}}
  \newunicodechar{≠}{\ensuremath{\neq}}
  \newunicodechar{≈}{\ensuremath{\approx}}
  \newunicodechar{∞}{\ensuremath{\infty}}
  \newunicodechar{∑}{\ensuremath{\sum}}
  \newunicodechar{∏}{\ensuremath{\prod}}
  \newunicodechar{∗}{\ensuremath{\ast}}
  \newunicodechar{·}{\ensuremath{\cdot}}
  \newunicodechar{−}{\ensuremath{-}}
  \newunicodechar{×}{\ensuremath{\times}}
  \newunicodechar{÷}{\ensuremath{\div}}
  \newunicodechar{±}{\ensuremath{\pm}}
  \newunicodechar{✓}{\ensuremath{\checkmark}}
  \newunicodechar{✗}{\ensuremath{\times}}
  \newunicodechar{⚠}{\textbf{!}}
  \newunicodechar{→}{\ensuremath{\rightarrow}}
  \newunicodechar{←}{\ensuremath{\leftarrow}}
  \newunicodechar{⇒}{\ensuremath{\Rightarrow}}
  \usepackage{xurl}
\fi

% --- Math ---
\usepackage{amsmath,amssymb,amsfonts,mathtools}
\usepackage{bm}

% --- Tables ---
\usepackage{booktabs}
\usepackage{longtable}
\usepackage{multirow}
\usepackage{array}
\usepackage{tabularx}

% --- Figures ---
\usepackage{graphicx}
\usepackage{float}
\usepackage{caption}
\captionsetup{font=small,labelfont=bf}

% --- Misc text ---
\usepackage{microtype}
\usepackage{enumitem}
\setlist{nosep,leftmargin=1.5em}
\usepackage{textcomp}
\usepackage{xcolor}
\usepackage{ulem}\normalem

% --- Code block support (pandoc Shaded environment) ---
\usepackage{fancyvrb}
\usepackage{framed}
\definecolor{shadecolor}{RGB}{248,248,248}
\newenvironment{Shaded}{\begin{snugshade}}{\end{snugshade}}
\providecommand{\BuiltInTok}[1]{#1}
\providecommand{\CommentTok}[1]{\textit{#1}}
\providecommand{\KeywordTok}[1]{\textbf{#1}}
\providecommand{\NormalTok}[1]{#1}
\providecommand{\StringTok}[1]{\textit{#1}}
\providecommand{\OperatorTok}[1]{#1}
\providecommand{\DataTypeTok}[1]{#1}
\providecommand{\FunctionTok}[1]{#1}


% --- Links & references ---
\usepackage[hidelinks,breaklinks=true]{hyperref}
\usepackage{cleveref}

% --- Bibliography (APA7 via natbib + apalike) ---
% APJM/MIR accept APA7. natbib's apalike close approximation; for full APA7
% compliance use biblatex with apa style.
\usepackage[round,sort]{natbib}
\bibliographystyle{apalike}

% --- Pandoc compatibility shims ---
\providecommand{\tightlist}{\setlength{\itemsep}{0pt}\setlength{\parskip}{0pt}}
\providecommand{\passthrough}[1]{\lstset{mathescape=false}#1\lstset{mathescape=true}}

% --- Title page (filled in per paper) ---
\title{P8\_PacificSIDS\_VN}
\author{Author details withheld for blind review}
\date{2026-05-22}

\begin{document}

\maketitle

\hypertarget{guxe1nh-nux1eb7ng-quux1ed1c-tux1ebf-huxf3a-bux1eaft-buux1ed9c-bux1eb1ng-chux1ee9ng-tux1eeb-cuxe1c-quux1ed1c-ux111ux1ea3o-nhux1ecf-ux111ang-phuxe1t-triux1ec3n-ux1edf-thuxe1i-buxecnh-dux1b0ux1a1ng}{%
\section{Gánh Nặng Quốc Tế Hóa Bắt Buộc: Bằng Chứng từ Các Quốc Đảo Nhỏ
Đang Phát Triển ở Thái Bình
Dương}\label{guxe1nh-nux1eb7ng-quux1ed1c-tux1ebf-huxf3a-bux1eaft-buux1ed9c-bux1eb1ng-chux1ee9ng-tux1eeb-cuxe1c-quux1ed1c-ux111ux1ea3o-nhux1ecf-ux111ang-phuxe1t-triux1ec3n-ux1edf-thuxe1i-buxecnh-dux1b0ux1a1ng}}

\emph{Bản dịch tiếng Việt --- Nguồn: p8/p8\_pacific\_sids\_en\_clean.md}
\emph{Dịch thuật: Claude AI theo chuẩn văn phong học thuật CTU
(09b\_vn\_term\_glossary.md v1.4)} \emph{Ngày: 17/05/2026} \emph{Tạp chí
mục tiêu: World Development (Elsevier)}

\begin{center}\rule{0.5\linewidth}{0.5pt}\end{center}

\textbf{Tạp chí mục tiêu}: Tạp chí Phát triển Thế giới (World
Development)

\textbf{Từ khóa}: quan hệ quốc tế hóa--hiệu quả doanh nghiệp; Quốc đảo
nhỏ đang phát triển; quốc tế hóa bắt buộc; khoảng trống thể chế; kinh tế
Thái Bình Dương; năng suất lao động

\textbf{Phân loại JEL:} F23; O19; L25; D22; O53.

\textbf{Loại bài:} Bài báo nghiên cứu.

\begin{center}\rule{0.5\linewidth}{0.5pt}\end{center}

\hypertarget{tuxf3m-tux1eaft}{%
\subsection{Tóm Tắt}\label{tuxf3m-tux1eaft}}

Giả thuyết hình chữ U ngược giữa cường độ quốc tế hóa và hiệu quả doanh
nghiệp (I→P) dựa trên ba điều kiện tiền đề cấu trúc: thị trường nội địa
khả thi, chi phí thương mại có thể quản lý được, và hỗ trợ thể chế có
chức năng. Chúng tôi lập luận rằng khi cả ba điều kiện này đồng thời
vắng mặt, như ở các Quốc đảo nhỏ đang phát triển (Small Island
Developing States, SIDS) ở Thái Bình Dương, mô hình phi tuyến tính kỳ
vọng sụp đổ thành quan hệ âm đơn điệu mà chúng tôi gọi là \textbf{Gánh
nặng quốc tế hóa bắt buộc (Forced Internationalization Penalty ---
FIP)}.

Không giống như độ dốc giảm ở giai đoạn 1 trong lý thuyết quốc tế hóa ba
giai đoạn, FIP không phải là chi phí học hỏi tạm thời với điểm ngưỡng
cuối cùng. FIP là điều kiện cấu trúc bền vững phát sinh từ sự bắt buộc
phải xuất khẩu bất chấp những bất lợi hệ thống.

Nghiên cứu sử dụng dữ liệu Khảo sát Doanh nghiệp của Ngân hàng Thế giới
(World Bank Enterprise Survey, WBES) từ 1.469 doanh nghiệp thuộc chín
nền kinh tế SIDS Thái Bình Dương (2009--2025). Kết quả phát hiện rằng
cường độ quốc tế hóa có quan hệ âm và vững chắc với năng suất lao động
(β = −0,404, p = 0,032 trong mô hình hiệu ứng cố định quốc gia-năm; β =
−1,236, p \textless{} 0,001 khi không có hiệu ứng cố định quốc gia),
không có độ cong bậc hai có ý nghĩa thống kê.

Năng lực công nghệ và mức độ áp dụng số, cả hai đều là biến điều tiết có
ý nghĩa trong bối cảnh châu Á đại lục, không có ý nghĩa thống kê trong
môi trường SIDS. Kết quả này phù hợp với cơ chế FIP: khi các điều kiện
tiền đề cấu trúc vắng mặt, năng lực cấp doanh nghiệp không thể chuyển
hướng quan hệ I→P sang vùng dương.

Các phát hiện này tái định khung giả thuyết hình chữ U ngược từ một quy
luật thực nghiệm phổ quát thành một quan hệ có điều kiện, xác định các
điều kiện biên cho lý thuyết điều tiết thể chế, và có hàm ý trực tiếp
cho Chương trình nghị sự Antigua và Barbuda dành cho SIDS.

\begin{center}\rule{0.5\linewidth}{0.5pt}\end{center}

\hypertarget{ux111iux1ec3m-nux1ed5i-bux1eadt-highlights}{%
\subsection{Điểm Nổi Bật
(Highlights)}\label{ux111iux1ec3m-nux1ed5i-bux1eadt-highlights}}

\begin{itemize}
\tightlist
\item
  FIP được định nghĩa là quan hệ âm đơn điệu của I→P khi ba điều kiện
  tiền đề cấu trúc đồng thời vắng mặt.
\item
  Bằng chứng từ 1.469 doanh nghiệp thuộc 9 nền kinh tế SIDS Thái Bình
  Dương (2009--2025).
\item
  β(FSTS\_c) = −0,404*, âm và có ý nghĩa thống kê; Lind--Mehlum xác nhận
  quan hệ âm đơn điệu.
\item
  Không tìm thấy điểm ngưỡng trong phạm vi quan sát, khác biệt hoàn toàn
  với hình chữ U ngược thông thường.
\item
  Năng lực công nghệ (TCI) và áp dụng số (DAI) không điều tiết được quan
  hệ I→P trong môi trường SIDS.
\item
  Hàm ý chính sách: chiến lược thúc đẩy xuất khẩu cần đi kèm can thiệp
  cấu trúc giải quyết chi phí thương mại, thị trường nội địa, và phát
  triển thể chế.
\end{itemize}

\begin{center}\rule{0.5\linewidth}{0.5pt}\end{center}

\hypertarget{1-giux1edbi-thiux1ec7u}{%
\subsection{1. Giới Thiệu}\label{1-giux1edbi-thiux1ec7u}}

\hypertarget{11-bux1ed1i-cux1ea3nh-vuxe0-vux1ea5n-ux111ux1ec1-nghiuxean-cux1ee9u}{%
\subsubsection{1.1 Bối Cảnh và Vấn Đề Nghiên
Cứu}\label{11-bux1ed1i-cux1ea3nh-vuxe0-vux1ea5n-ux111ux1ec1-nghiuxean-cux1ee9u}}

Quan hệ giữa quốc tế hóa và hiệu quả doanh nghiệp (I→P) là một trong
những quy luật thực nghiệm gây tranh luận nhiều nhất trong nghiên cứu
kinh doanh quốc tế (international business, IB). Các phân tích tổng hợp
nhất quán báo cáo tác động dương tích cực trung bình (Kirca et al.,
2011; Bausch \& Krist, 2007). Tuy nhiên, phương sai giữa các nghiên cứu
rất lớn và không đồng nhất; đây là lời mời để khảo sát các \emph{điều
kiện} mà dưới đó quốc tế hóa giúp ích hoặc gây hại cho hiệu quả doanh
nghiệp (Verbeke \& Brugman, 2018).

Giải pháp lý thuyết chủ đạo giả định dạng hàm hình chữ U ngược (Lu \&
Beamish, 2004). Quốc tế hóa ban đầu cải thiện hiệu quả thông qua kinh tế
quy mô, học hỏi, và đa dạng hóa thị trường. Sau đó, vượt qua một ngưỡng
nhất định, chi phí phối hợp và độ phức tạp tổ chức làm xói mòn những lợi
ích này. Contractor, Kundu và Hsu (2003) mở rộng điều này thành đường
cong S ba giai đoạn, trong đó ngay cả độ dốc giảm ban đầu ở giai đoạn 1
cũng chỉ là tạm thời và cuối cùng nhường chỗ cho lợi ích ròng.

Cả hai khung lý thuyết đều chia sẻ một giả định nền tảng chung: môi
trường thể chế, trong khi có thể điều tiết \emph{cường độ} của quan hệ
I→P, không làm thay đổi hình dạng cơ bản hay dấu hiệu của quan hệ đó.

Bài báo này thách thức giả định đó. Chúng tôi lập luận rằng một số điều
kiện cấu trúc, đặc trưng của các Quốc đảo nhỏ đang phát triển (SIDS),
loại bỏ các điều kiện tiền đề mà logic hình chữ U ngược phụ thuộc vào.
Khi đó, mô hình phi tuyến tính kỳ vọng được thay thế bằng quan hệ âm đơn
điệu. Chúng tôi gọi hiện tượng này là \textbf{Gánh nặng quốc tế hóa bắt
buộc (Forced Internationalization Penalty, FIP)}: tác động âm bền vững
của cường độ quốc tế hóa lên hiệu quả doanh nghiệp trong các bối cảnh mà
ba điều kiện tiền đề cấu trúc đồng thời vắng mặt.

\hypertarget{12-bux1ed1i-cux1ea3nh-sids-vuxe0-luxfd-do-chux1ecdn}{%
\subsubsection{1.2 Bối Cảnh SIDS và Lý Do
Chọn}\label{12-bux1ed1i-cux1ea3nh-sids-vuxe0-luxfd-do-chux1ecdn}}

SIDS đại diện cho bối cảnh có tính cực đoan về lý thuyết nhưng lại thiếu
hụt về mặt nghiên cứu thực nghiệm trong lĩnh vực I→P. Các nền kinh tế
này được đặc trưng bởi ba điểm. Thứ nhất, thị trường nội địa cực nhỏ
(Tonga: 104.000 dân; Kiribati: 121.000 dân). Thứ hai, cô lập địa lý cực
đoan làm khuếch đại chi phí thương mại lên 30--210\% (Winters \&
Martins, 2004). Thứ ba, khoảng trống thể chế toàn diện hạn chế khả năng
chiếm đoạt lợi nhuận của doanh nghiệp từ các hoạt động quốc tế
(Briguglio, 1995; Khanna \& Palepu, 2010).

Theo dõi hệ thống gần đây trên 121.249 điểm dữ liệu phát triển xác nhận
rằng các điều kiện này đã ngày càng trầm trọng hơn. Tiến độ giai đoạn
2020--2025 là chậm nhất từ trước đến nay trên 15 trong số 26 tiêu chuẩn
được theo dõi, với SIDS và châu Phi cận Sahara bị ảnh hưởng không cân
xứng (Pirlea et al., 2026). Phân tích đa chỉ số độc lập trên sáu chiều
phát triển (nghèo đói, tuổi thọ, giáo dục, bình đẳng giới, khí hậu, và
điện) cho thấy thế giới đang phát triển với tốc độ chậm nhất trong 75
năm qua. Có 54 nền kinh tế còn hơn một thế kỷ nữa mới đạt tiêu chuẩn
nước phát triển theo quỹ đạo hiện tại, và SIDS chiếm tỷ lệ không cân
xứng trong nhóm này (Mahler, Serajuddin, Wadhwa, \& Yonzan, 2026).

Điều quan trọng là các doanh nghiệp ở SIDS không \emph{chọn} quốc tế hóa
như một chiến lược tăng trưởng; đối với nhiều doanh nghiệp, xuất khẩu là
mệnh lệnh cấu trúc được quy định bởi sự bất khả thi của việc đạt quy mô
hiệu quả tối thiểu trong nước. Sự bắt buộc này, quốc tế hóa vì sinh tồn,
thay đổi cơ bản logic của quan hệ I→P.

\hypertarget{13-ux111uxf3ng-guxf3p-cux1ee7a-buxe0i-buxe1o}{%
\subsubsection{1.3 Đóng Góp của Bài
Báo}\label{13-ux111uxf3ng-guxf3p-cux1ee7a-buxe0i-buxe1o}}

Bài báo này có ba đóng góp vào văn liệu IB. Thứ nhất, chúng tôi giới
thiệu và định nghĩa lý thuyết FIP như một điều kiện biên của mô hình
hình chữ U ngược. Khái niệm FIP được phân biệt với các khái niệm liền kề
bao gồm chi phí giai đoạn 1 trong lý thuyết ba giai đoạn (Contractor et
al., 2003), Trách nhiệm ngoại lai (Liability of Foreignness, Zaheer,
1995), và khởi nghiệp bắt buộc (necessity entrepreneurship, Sarasvathy
et al., 2014).

Thứ hai, chúng tôi cung cấp bằng chứng thực nghiệm từ 1.469 doanh nghiệp
thuộc chín nền kinh tế SIDS Thái Bình Dương bằng dữ liệu WBES, phân tích
I→P cấp doanh nghiệp toàn diện nhất tại SIDS Thái Bình Dương cho đến
nay.

Thứ ba, chúng tôi ghi lại rằng năng lực cấp doanh nghiệp (Chỉ số năng
lực công nghệ, TCI, Chỉ số áp dụng số, DAI) không có khả năng vượt qua
FIP, chuyển hướng cuộc thảo luận lý thuyết từ "điều tiết thể chế của
quan hệ I→P" (Marano, Tallman, \& Teece, 2016) sang "quyết định thể chế
về dấu của quan hệ I→P."

Các phát hiện có hàm ý chính sách trực tiếp theo Chương trình nghị sự
Antigua và Barbuda dành cho SIDS (United Nations General Assembly,
2024), cho thấy các chiến lược dựa trên tăng trưởng xuất khẩu có thể
không phù hợp một cách hệ thống với thực tế cấu trúc của các vi kinh tế
Thái Bình Dương.

\begin{center}\rule{0.5\linewidth}{0.5pt}\end{center}

\hypertarget{2-khung-luxfd-thuyux1ebft-vuxe0-giux1ea3-thuyux1ebft}{%
\subsection{2. Khung Lý Thuyết và Giả
Thuyết}\label{2-khung-luxfd-thuyux1ebft-vuxe0-giux1ea3-thuyux1ebft}}

\hypertarget{21-muxf4-huxecnh-huxecnh-chux1eef-u-ngux1b0ux1ee3c-vuxe0-cuxe1c-ux111iux1ec1u-kiux1ec7n-tiux1ec1n-ux111ux1ec1-cux1ea5u-truxfac}{%
\subsubsection{2.1 Mô Hình Hình Chữ U Ngược và Các Điều Kiện Tiền Đề Cấu
Trúc}\label{21-muxf4-huxecnh-huxecnh-chux1eef-u-ngux1b0ux1ee3c-vuxe0-cuxe1c-ux111iux1ec1u-kiux1ec7n-tiux1ec1n-ux111ux1ec1-cux1ea5u-truxfac}}

Giả thuyết hình chữ U ngược trong nghiên cứu I→P xuất phát từ logic chi
phí-lợi ích. Khi doanh nghiệp mở rộng sang thị trường nước ngoài, các
tác động dương ban đầu (kinh tế quy mô, thu nhận kiến thức, đa dạng hóa
thị trường) chiếm ưu thế. Vượt qua một ngưỡng nhất định, độ phức tạp tổ
chức ngày càng tăng, cùng chi phí phối hợp và trách nhiệm ngoại lai đảo
ngược quỹ đạo hiệu quả (Lu \& Beamish, 2004; Hitt, Hoskisson, \& Kim,
1997). Điểm ngưỡng, thường được ước tính trong khoảng 30\% đến 70\%
doanh số nước ngoài trên tổng doanh số (cường độ xuất khẩu, Foreign
Sales to Total Sales, FSTS), đại diện cho điểm tại đó chi phí biên bằng
lợi ích biên của quốc tế hóa tiếp theo.

Logic này dựa trên, một cách ngầm định, ba điều kiện tiền đề cấu trúc
hiếm khi được phát biểu nhưng thường được thỏa mãn trong các mẫu nghiên
cứu IB chính lưu:

\textbf{Điều kiện tiền đề 1: Thị trường nội địa khả thi.} Logic hình chữ
U ngược giả định rằng doanh nghiệp có căn cứ nội địa hoạt động cung cấp
doanh thu, tài nguyên, và học hỏi tổ chức trước và trong quá trình quốc
tế hóa (Johanson \& Vahlne, 1977; 2009). Khi thị trường nội địa không
khả thi, quá nhỏ để hỗ trợ quy mô hiệu quả tối thiểu cho ngay cả các
hoạt động cơ bản, doanh nghiệp bị ép buộc về mặt cấu trúc phải xuất khẩu
bất kể sự sẵn sàng chiến lược hay năng lực. Giả định "lựa chọn xuất
khẩu" làm cơ sở cho hầu hết các mô hình I→P bị phá vỡ.

\textbf{Điều kiện tiền đề 2: Chi phí thương mại có thể quản lý được.} Mô
hình tiêu chuẩn giả định rằng, ít nhất ban đầu, doanh thu xuất khẩu vượt
chi phí xuất khẩu, cho phép một giai đoạn lợi nhuận quốc tế hóa ròng
dương. Ở các nền kinh tế bị cô lập về địa lý, chi phí thương mại có thể
cực cao đến mức điều kiện này bị vi phạm ngay cả ở cường độ xuất khẩu
nhỏ. Winters và Martins (2004) ghi lại mức phí bù chi phí thương mại từ
30--210\% cho các nền kinh tế nhỏ xa xôi, phần lớn độc lập với hiệu quả
doanh nghiệp hay lựa chọn chiến lược.

\textbf{Điều kiện tiền đề 3: Hỗ trợ thể chế có chức năng.} Cả lý thuyết
thể chế (Scott, 1995; North, 1990) và văn liệu khoảng trống thể chế
(Khanna \& Palepu, 2010; Doh et al., 2017) đều coi chất lượng thể chế hỗ
trợ thị trường là biến điều tiết cường độ quan hệ I→P. Điều này giả định
rằng thể chế biến thiên trên một thang liên tục nơi tồn tại một mức hỗ
trợ chức năng nhất định ngay cả ở mức thấp nhất. Khoảng trống thể chế
toàn diện, thiếu vắng các trung gian tài chính, thực thi hợp đồng không
đáng tin cậy, hạ tầng thuận lợi thương mại hạn chế, không chỉ làm suy
yếu đường cong hình chữ U ngược; chúng có thể loại bỏ hoàn toàn cơ chế
tạo ra lợi nhuận dương từ quốc tế hóa.

\hypertarget{22-guxe1nh-nux1eb7ng-quux1ed1c-tux1ebf-huxf3a-bux1eaft-buux1ed9c-ux111ux1ecbnh-nghux129a-luxfd-thuyux1ebft}{%
\subsubsection{2.2 Gánh Nặng Quốc Tế Hóa Bắt Buộc: Định Nghĩa Lý
Thuyết}\label{22-guxe1nh-nux1eb7ng-quux1ed1c-tux1ebf-huxf3a-bux1eaft-buux1ed9c-ux111ux1ecbnh-nghux129a-luxfd-thuyux1ebft}}

Chúng tôi định nghĩa \textbf{Gánh nặng quốc tế hóa bắt buộc (Forced
Internationalization Penalty --- FIP)} là tác động âm đơn điệu của cường
độ quốc tế hóa lên hiệu quả doanh nghiệp trong các bối cảnh mà ba điều
kiện tiền đề cấu trúc nêu trên đồng thời vắng mặt. FIP khác với các khái
niệm liên quan ở bốn điểm:

\emph{So với chi phí giai đoạn 1 (Contractor et al., 2003)}: Chi phí
giai đoạn 1 là tạm thời, chúng phản ánh chi phí đầu tư thiết lập và học
hỏi ban đầu mà cuối cùng tạo ra lợi ích ròng một khi vượt qua điểm
ngưỡng. FIP không có điểm ngưỡng lý thuyết trong phạm vi quan sát của
FSTS, vì ràng buộc mang tính cấu trúc (thị trường nội địa không khả thi,
chi phí thương mại cực cao, khoảng trống thể chế) chứ không phải chuyển
tiếp.

\emph{So với Trách nhiệm ngoại lai (Zaheer, 1995)}: Trách nhiệm ngoại
lai biểu thị chi phí bổ sung mà các doanh nghiệp nước ngoài phải chịu
khi hoạt động ở thị trường nước chủ nhà so với các đối thủ trong nước.
FIP hoạt động ở phía \emph{thị trường nội địa} của phương trình, đây là
gánh nặng mà các doanh nghiệp SIDS phải chịu khi họ xuất khẩu, không
phải gánh nặng mà các tập đoàn đa quốc gia phải chịu khi thâm nhập thị
trường SIDS.

\emph{So với khởi nghiệp bắt buộc (Sarasvathy et al., 2014)}: Các nhà
khởi nghiệp bắt buộc trong các bối cảnh phát triển thường đạt được ngang
bằng hiệu quả với các nhà khởi nghiệp có định hướng cơ hội sau đủ thời
gian và hỗ trợ thể chế. FIP đặc trưng cho các bối cảnh nơi hỗ trợ thể
chế vắng mặt một cách hệ thống, làm cho con đường quốc tế hóa do bắt
buộc luôn kém lợi thế.

\emph{So với tính dễ bị tổn thương kinh tế SIDS (Briguglio, 1995)}: Chỉ
số dễ bị tổn thương của Briguglio nắm bắt mức độ tiếp xúc kinh tế vĩ mô
với các cú sốc bên ngoài. FIP là một cấu trúc kinh tế vi mô, cấp doanh
nghiệp: nó xác định cơ chế hiệu quả thông qua đó tính dễ bị tổn thương
cấu trúc vĩ mô chuyển hóa thành kết quả cấp doanh nghiệp.

\hypertarget{23-bux1ed1i-cux1ea3nh-sids-vuxe0-quux1ed1c-tux1ebf-huxf3a-vuxec-sinh-tux1ed3n}{%
\subsubsection{2.3 Bối Cảnh SIDS và Quốc Tế Hóa Vì Sinh
Tồn}\label{23-bux1ed1i-cux1ea3nh-sids-vuxe0-quux1ed1c-tux1ebf-huxf3a-vuxec-sinh-tux1ed3n}}

Các nền kinh tế SIDS Thái Bình Dương thể hiện cả ba điều kiện tiền đề bị
vi phạm ở dạng cực đoan. Với dân số dao động từ 94.000 (Micronesia) đến
10 triệu (Papua New Guinea, nền kinh tế lớn nhất trong mẫu của chúng
tôi), thị trường nội địa không thể duy trì các doanh nghiệp sản xuất hay
dịch vụ ở quy mô hiệu quả tối thiểu. Read (2004) chứng minh rằng các nền
kinh tế đảo nhỏ nhất thực tế không có kiến trúc lựa chọn để kiêng xuất
khẩu, các doanh nghiệp hoặc phải xuất khẩu hoặc phải rút khỏi thị
trường. Chúng tôi gọi mô hình này là \emph{quốc tế hóa vì sinh tồn}:
không phải là mở rộng chiến lược mà là sự bắt buộc cấu trúc.

Trong quốc tế hóa vì sinh tồn, FSTS, thước đo tiêu chuẩn của cường độ
quốc tế hóa, không còn lập chỉ số cam kết tự nguyện vào thị trường nước
ngoài; nó ủy quyền cho \emph{mức độ không đủ năng lực của thị trường nội
địa}. Các doanh nghiệp có FSTS cao hơn không nhất thiết là những doanh
nghiệp có năng lực quốc tế vượt trội; họ có thể chỉ đơn giản là đối mặt
với các ràng buộc thị trường nội địa cấp tính hơn. Cách tái diễn giải
FSTS trong bối cảnh SIDS này tạo ra một dự báo thực nghiệm mới:

\textbf{Giả thuyết 1 (Giả thuyết FIP)}: Trong SIDS Thái Bình Dương,
cường độ quốc tế hóa có quan hệ âm với hiệu quả doanh nghiệp mà không có
điểm ngưỡng bậc hai có ý nghĩa thống kê, cho thấy Gánh nặng quốc tế hóa
bắt buộc chứ không phải quan hệ hình chữ U ngược.

\hypertarget{24-vai-truxf2-cux1ee7a-nux103ng-lux1ef1c-cux1ea5p-doanh-nghiux1ec7p}{%
\subsubsection{2.4 Vai Trò của Năng Lực Cấp Doanh
Nghiệp}\label{24-vai-truxf2-cux1ee7a-nux103ng-lux1ef1c-cux1ea5p-doanh-nghiux1ec7p}}

Trong bối cảnh châu Á đại lục và các bối cảnh thể chế khác, năng lực cấp
doanh nghiệp điều tiết quan hệ I→P. Chỉ số năng lực công nghệ
(Technological Capability Index, TCI) (Lall, 1992; Goedhuys \&
Sleuwaegen, 2013) và Chỉ số áp dụng số (Digital Adoption Index, DAI)
(Teece, 2007) được chứng minh là dịch chuyển đường cong I→P lên phía
trên, cho phép doanh nghiệp đạt hiệu quả tốt hơn ở bất kỳ mức độ quốc tế
hóa nào, và thay đổi điểm ngưỡng. Cơ chế lý thuyết là các năng lực làm
giảm chi phí phức tạp tổ chức ở phần giảm dần của hình chữ U ngược.

Trong bối cảnh FIP, cơ chế điều tiết này có thể bị suy giảm. Nếu quan hệ
I→P âm xuất phát từ các điều kiện cấu trúc (thị trường nội địa không khả
thi, chi phí thương mại cực cao, khoảng trống thể chế) thay vì từ sự
phức tạp tổ chức, thì các năng lực cấp doanh nghiệp giải quyết sự phức
tạp tổ chức sẽ không chuyển hướng quan hệ này. Một doanh nghiệp SIDS với
năng lực công nghệ vượt trội vẫn phải đối mặt với mức khuếch đại chi phí
thương mại 30--210\%; một doanh nghiệp SIDS có trang web vẫn phải đối
mặt với sự vắng mặt của các trung gian tài chính có chức năng cho tài
trợ thương mại. Cơ chế cấu trúc tạo ra FIP hoạt động ở mức độ mà các
năng lực cấp doanh nghiệp không thể dễ dàng bù đắp.

\textbf{Giả thuyết 2 (Giả thuyết null năng lực)}: Trong SIDS Thái Bình
Dương, Chỉ số năng lực công nghệ (TCI) và Chỉ số áp dụng số (DAI) không
điều tiết có ý nghĩa thống kê quan hệ quốc tế hóa--hiệu quả doanh
nghiệp.

\emph{Lưu ý}: H2 được phát biểu là giả thuyết với kỳ vọng null dựa trên
cơ chế lý thuyết FIP, không phải là lập luận về năng lực thống kê. Cách
đặt vấn đề này theo hướng dẫn phương pháp của Haans, Pieters, và He
(2016) về việc đặt giả thuyết cho kết quả null vì lý do có cơ sở lý
thuyết.

\begin{center}\rule{0.5\linewidth}{0.5pt}\end{center}

\hypertarget{3-dux1eef-liux1ec7u-vuxe0-phux1b0ux1a1ng-phuxe1p}{%
\subsection{3. Dữ Liệu và Phương
Pháp}\label{3-dux1eef-liux1ec7u-vuxe0-phux1b0ux1a1ng-phuxe1p}}

\hypertarget{31-nguux1ed3n-dux1eef-liux1ec7u}{%
\subsubsection{3.1 Nguồn Dữ
Liệu}\label{31-nguux1ed3n-dux1eef-liux1ec7u}}

Chúng tôi sử dụng dữ liệu tổng hợp từ Khảo sát Doanh nghiệp Ngân hàng
Thế giới (World Bank Enterprise Survey --- WBES), lấy từ các đợt khảo
sát được thực hiện ở các nền kinh tế SIDS Thái Bình Dương trong giai
đoạn 2009--2025. WBES sử dụng lấy mẫu ngẫu nhiên phân tầng với thiết kế
xác suất tỷ lệ với quy mô trên các tầng quy mô doanh nghiệp (5--19,
20--99, 100+ nhân viên) và các ngành (sản xuất, dịch vụ, bán lẻ).

Chúng tôi đưa vào tất cả các quan sát doanh nghiệp-năm từ các nền kinh
tế đáp ứng phân loại SIDS: Fiji, Kiribati, Papua New Guinea, Samoa,
Solomon Islands, Timor-Leste, Tonga, Vanuatu, và Comoros (chín nền kinh
tế). Mẫu phân tích, các quan sát có giá trị không thiếu về năng suất lao
động, cường độ xuất khẩu, và tuổi doanh nghiệp, bao gồm \textbf{1.469
doanh nghiệp} trên nhiều đợt quốc gia-năm.

\hypertarget{32-ux111o-lux1b0ux1eddng-biux1ebfn}{%
\subsubsection{3.2 Đo Lường
Biến}\label{32-ux111o-lux1b0ux1eddng-biux1ebfn}}

\textbf{Biến phụ thuộc.} Năng suất lao động log (ln\_LP) được xây dựng
là logarithm tự nhiên của doanh thu bán hàng trên mỗi lao động thường
xuyên, tính bằng USD đã điều chỉnh theo sức mua tương đương (PPP) không
đổi: ln(d2/l1). Thước đo này được sử dụng rộng rãi trong các nghiên cứu
năng suất cấp doanh nghiệp (Goedhuys \& Sleuwaegen, 2013; World Bank,
2020). Trong số 1.469 quan sát phân tích, 187 quan sát (12,7\%) là các
nhà xuất khẩu với FSTS \textgreater{} 0.

\textbf{Cường độ quốc tế hóa.} Theo phương pháp của Haans et al. (2016)
và Contractor et al. (2003), chúng tôi xây dựng FSTS (cường độ xuất khẩu
--- Foreign Sales to Total Sales) từ mục WBES d3c (tỷ lệ phần trăm doanh
số xuất khẩu), chia cho 100. Đối với các mô hình chính, FSTS được cân
bằng theo trung bình đợt khảo sát: FSTS\_c = FSTS −
trung\_bình\_đợt(FSTS), cho trung bình FSTS là 0,060 (độ lệch chuẩn
không báo cáo do mẫu con nhà xuất khẩu nhỏ). Cân bằng trung bình triệt
tiêu sự khác biệt mức độ giữa các đợt trong khi bảo toàn biến thiên
trong đợt, theo thực hành tiêu chuẩn cho các nghiên cứu kiểm tra độ cong
FSTS (Haans et al., 2016).

\textbf{Chỉ số năng lực công nghệ (Technological Capability Index ---
TCI\_z).} Theo Lall (1992) và Goedhuys và Sleuwaegen (2013), TCI được
xây dựng từ hai mục WBES: b8 (chứng nhận chất lượng được công nhận quốc
tế, nhị phân) và e6 (cấp phép công nghệ nước ngoài, nhị phân). TCI\_z là
giá trị trung bình được chuẩn hóa z của hai chỉ số này.

\textbf{Chỉ số áp dụng số (Digital Adoption Index --- DAI\_z).} DAI được
đo bằng mục WBES c22b (doanh nghiệp có trang web hoạt động, nhị phân).
Đây là proxy Tầng 1 nắm bắt \emph{sự hiện diện trang web nền tảng}, bước
đầu tiên của sự hiện diện số. Thước đo này phù hợp với bối cảnh SIDS nơi
hạ tầng số tiên tiến còn hạn chế, chỉ 5\% dân số các nước thu nhập thấp
có kỹ năng số cơ bản vào năm 2023, và các nền kinh tế thu nhập cao chiếm
77\% công suất trung tâm dữ liệu toàn cầu (Mahler, Wang, \& Weber,
2026). DAI\_z là giá trị chuẩn hóa z. Chúng tôi không nhận định rằng
mình đang đo năng lực số động, vì điều này đòi hỏi các chỉ số đa tầng
không có sẵn trong WBES cho mẫu này.

\textbf{Biến kiểm soát.} Quy mô doanh nghiệp là logarithm tự nhiên của
số lao động thường xuyên (ln\_size, từ l1). Tuổi doanh nghiệp là năm
khảo sát trừ năm thành lập (b5). Tỷ lệ sở hữu nước ngoài
(foreign\_own\_pct) là tỷ lệ phần trăm vốn chủ sở hữu doanh nghiệp do
các thực thể nước ngoài nắm giữ (b2b). Đây là các biến kiểm soát tiêu
chuẩn trong nghiên cứu I→P cấp doanh nghiệp (Hitt et al., 1997; Lu \&
Beamish, 2004).

\emph{Bảng 0: Định nghĩa biến và cách xây dựng từ WBES.}

\begin{longtable}[]{@{}llll@{}}
\toprule\noalign{}
Biến & Mục WBES & Cách xây dựng & Vai trò trong mô hình \\
\midrule\noalign{}
\endhead
\bottomrule\noalign{}
\endlastfoot
ln\_LP & d2, l1 & ln(doanh thu hàng năm PPP / lao động thường xuyên);
điều chỉnh PPP sử dụng hệ số giảm phát World Bank ICP theo năm khảo sát
& Biến phụ thuộc --- log năng suất lao động (Goedhuys \& Sleuwaegen,
2013) \\
FSTS & d3c & d3c / 100; tỷ lệ tổng doanh số từ xuất khẩu & Cường độ quốc
tế hóa --- thô (Haans et al., 2016) \\
FSTS\_c & d3c & FSTS − trung\_bình\_đợt(FSTS); cân bằng trung bình theo
đợt & Cường độ quốc tế hóa --- đã cân bằng (Contractor et al., 2003) \\
FSTS\_c² & d3c & FSTS\_c²; kiểm định độ cong (H1) & Hạng tử bậc hai; β₂
\textless{} 0 nếu hình chữ U ngược, KÝ HIỆU KHÔNG ĐÁG nếu đơn điệu FIP
(Lind \& Mehlum, 2010) \\
TCI\_z & b8, e6 & chuẩn hóa-z(trung\_bình(b8\_nhị\_phân,
e6\_nhị\_phân)); chứng nhận chất lượng + giấy phép công nghệ nước ngoài
& Năng lực công nghệ --- proxy Tầng 1; kiểm định trực tiếp + điều tiết
(H2; Lall, 1992) \\
DAI\_z & c22b & chuẩn hóa-z(c22b\_nhị\_phân); sự hiện diện trang web
hoạt động --- chỉ Tầng 1 & Áp dụng số --- proxy trang web nền tảng;
KHÔNG phải năng lực số động (Banalieva \& Dhanaraj, 2019; H2) \\
ln\_size & l1 & ln(lao động thường xuyên) & Biến kiểm soát --- quy mô
doanh nghiệp (Hitt et al., 1997) \\
firm\_age & b5 & năm\_khảo\_sát − b5 & Biến kiểm soát --- tuổi doanh
nghiệp (Lu \& Beamish, 2004) \\
foreign\_own\_pct & b2b & \% vốn chủ sở hữu do thực thể nước ngoài nắm
giữ & Biến kiểm soát --- sở hữu nước ngoài (Marano et al., 2016) \\
δ\_c & a4a/a4b & Hiệu ứng cố định quốc gia (9 nền kinh tế SIDS) & Hấp
thụ tính không đồng nhất cấp quốc gia bất biến theo thời gian: chất
lượng thể chế, sự xa xôi, chi phí thương mại \\
λ\_t & đợt & Hiệu ứng cố định đợt-năm (2009, 2015, 2016, 2018, 2019,
2022, 2023, 2024, 2025) & Hấp thụ xu hướng thời gian chung \\
\end{longtable}

\emph{Ghi chú.} b8 và e6 được mã hóa lại từ mã phản hồi WBES (1 = có, 2
= không) sang nhị phân (1/0). c22b được mã hóa lại tương tự. TCI\_z và
DAI\_z được chuẩn hóa trên toàn bộ mẫu tổng hợp. Chuyển đổi PPP sử dụng
hệ số giảm phát World Bank ICP khớp với năm khảo sát. FSTS\_c được cân
bằng trong từng đợt quốc gia-năm để các ước lượng điểm ngưỡng không bị
nhiễu bởi sự dịch chuyển cơ cấu giữa các đợt (Haans et al., 2016). DAI
chỉ nắm bắt sự hiện diện trang web tĩnh Tầng 1; các chỉ số Tầng 2 (thanh
toán điện tử) và Tầng 3 (điện toán đám mây/ERP) không có sẵn cho các đợt
WBES này.

\hypertarget{33-chuux1ed7i-muxf4-huxecnh}{%
\subsubsection{3.3 Chuỗi Mô Hình}\label{33-chuux1ed7i-muxf4-huxecnh}}

Chúng tôi ước tính bốn họ mô hình để kiểm định H1 và H2:

\textbf{M0 (cơ sở kiểm soát)}: ln\_LP\_it = α + γ₁ ln\_size\_it + γ₂
firm\_age\_it + γ₃ foreign\_own\_pct\_it + δ\_c + λ\_t + ε\_it

\textbf{M1 (FSTS tuyến tính, kiểm định FIP)}: ln\_LP\_it = α + β₁
FSTS\_c\_it + γ·X\_it + δ\_c + λ\_t + ε\_it

\textbf{M2 (FSTS bậc hai, kiểm định độ cong)}: ln\_LP\_it = α + β₁
FSTS\_c\_it + β₂ FSTS\_c²\_it + γ·X\_it + δ\_c + λ\_t + ε\_it

\textbf{M3 (năng lực, kiểm định null điều tiết)}: ln\_LP\_it = α + β₁
FSTS\_c\_it + β₂ FSTS\_c²\_it + β₃ TCI\_z\_it + β₄ DAI\_z\_it + γ·X\_it
+ δ\_c + λ\_t + ε\_it

trong đó δ\_c biểu thị hiệu ứng cố định quốc gia, λ\_t hiệu ứng cố định
đợt-năm, và X vector biến kiểm soát. Hiệu ứng cố định quốc gia hấp thụ
tính không đồng nhất cấp quốc gia bất biến theo thời gian (chất lượng
thể chế, chi phí thương mại, sự xa xôi địa lý) mà nếu không sẽ làm nhiễu
hệ số FSTS.

Chúng tôi cũng ước tính hai đặc tả kiểm định vững bổ sung: M\_yearFE
(chỉ hiệu ứng cố định năm, không có hiệu ứng cố định quốc gia) và
M\_bivariate (không có biến kiểm soát), để đánh giá độ nhạy của ước
lượng FIP theo các giả định mô hình hóa khác nhau.

Tất cả các hồi quy được ước tính bằng OLS (Ordinary Least Squares, bình
phương tối thiểu thông thường) với sai số chuẩn vững với phương sai
không đồng nhất (HC1). Phân tích được thực hiện trong R; tập lệnh nhân
rộng đầy đủ, bảng hệ số, và thống kê tóm tắt được lưu trữ trong
\texttt{p8/replication/}.

\hypertarget{34-kiux1ec3m-ux111ux1ecbnh-vux1eefng}{%
\subsubsection{3.4 Kiểm Định
Vững}\label{34-kiux1ec3m-ux111ux1ecbnh-vux1eefng}}

Để đảm bảo kết quả FIP không phụ thuộc vào một đặc tả mô hình đơn lẻ,
chúng tôi thực hiện ba kiểm định vững bổ sung:

\begin{enumerate}
\def\labelenumi{\arabic{enumi}.}
\tightlist
\item
  \textbf{M\_yearFE}: Loại bỏ hiệu ứng cố định quốc gia để đánh giá xem
  FIP có duy trì hay không khi khai thác biến thiên giữa các quốc gia.
\item
  \textbf{M\_bivariate}: Hồi quy song biến không có biến kiểm soát, kiểm
  tra xem mô hình đơn giản nhất có xác nhận chiều âm hay không.
\item
  \textbf{Mẫu con chỉ nhà xuất khẩu (N = 187)}: Giới hạn phân tích chỉ
  cho các doanh nghiệp có FSTS \textgreater{} 0 để loại trừ giải thích
  lựa chọn bất lợi.
\end{enumerate}

\begin{center}\rule{0.5\linewidth}{0.5pt}\end{center}

\hypertarget{4-kux1ebft-quux1ea3}{%
\subsection{4. Kết Quả}\label{4-kux1ebft-quux1ea3}}

\hypertarget{41-muxf4-tux1ea3-mux1eabu-vuxe0-bux1eb1ng-chux1ee9ng-sux1a1-bux1ed9}{%
\subsubsection{4.1 Mô Tả Mẫu và Bằng Chứng Sơ
Bộ}\label{41-muxf4-tux1ea3-mux1eabu-vuxe0-bux1eb1ng-chux1ee9ng-sux1a1-bux1ed9}}

Mẫu phân tích bao gồm 1.469 doanh nghiệp thuộc chín nền kinh tế SIDS
Thái Bình Dương. Bảng 1 báo cáo thống kê tóm tắt cấp quốc gia. Tỷ lệ đại
diện nhà xuất khẩu thấp trên tất cả các nền kinh tế (tổng thể 12,7\%),
phù hợp với ràng buộc cấu trúc SIDS: hầu hết các doanh nghiệp hoạt động
chủ yếu ở thị trường nội địa bất chấp các hạn chế của nó. FSTS trung
bình trên toàn mẫu là 0,060, phản ánh sự tham gia xuất khẩu thấp nhưng
không bằng không của các doanh nghiệp SIDS.

Hiệu ứng cố định quốc gia hấp thụ phương sai đáng kể. R² tăng từ 0,511
(chỉ hiệu ứng cố định năm) lên 0,799 (hiệu ứng cố định quốc gia + năm)
khi thêm các chỉ số quốc gia. Kết quả này xác nhận rằng 95\% phương sai
năng suất được giải thích là giữa các quốc gia, một đặc điểm phù hợp với
tính không đồng nhất cực cao giữa các quốc gia về chất lượng thể chế, sự
xa xôi địa lý, và cấu trúc kinh tế.

\hypertarget{42-kux1ebft-quux1ea3-chuxednh-guxe1nh-nux1eb7ng-quux1ed1c-tux1ebf-huxf3a-bux1eaft-buux1ed9c-h1}{%
\subsubsection{4.2 Kết Quả Chính: Gánh Nặng Quốc Tế Hóa Bắt Buộc
(H1)}\label{42-kux1ebft-quux1ea3-chuxednh-guxe1nh-nux1eb7ng-quux1ed1c-tux1ebf-huxf3a-bux1eaft-buux1ed9c-h1}}

Bảng 2 báo cáo các kết quả hồi quy chính. \textbf{M1}, đặc tả FSTS tuyến
tính với hiệu ứng cố định quốc gia và năm, cho kết quả:

β(FSTS\_c) = \textbf{−0,404}, SE = 0,188, t = −2,152, p = \textbf{0,032}

Hệ số này có ý nghĩa thống kê ở mức 5\% thông thường và có ý nghĩa thực
chất: một đơn vị tăng trong FSTS (di chuyển từ không xuất khẩu sang định
hướng xuất khẩu hoàn toàn) gắn liền với năng suất lao động thấp hơn
40,4\%. Điều này ủng hộ \textbf{Giả thuyết 1}: cường độ quốc tế hóa có
quan hệ âm với hiệu quả doanh nghiệp trong SIDS Thái Bình Dương.

Cơ sở M0 (chỉ biến kiểm soát, không có FSTS) cho R² là 0,7992; thêm
FSTS\_c trong M1 tăng R² một chút lên 0,8001, cho thấy rằng, có điều
kiện trên hiệu ứng cố định quốc gia và năm, FSTS giải thích phương sai
bổ sung hạn chế nhưng theo chiều âm nhất quán.

Kết quả biến kiểm soát trong M1 nhất quán với bằng chứng đã thiết lập.
Tuổi doanh nghiệp có quan hệ dương với năng suất (β = +0,009, p =
0,052). Sở hữu nước ngoài cho thấy quan hệ âm ở mức biên (β = −0,006, p
= 0,056), có thể phản ánh hiệu quả đầu tư nước ngoài thấp trong bối cảnh
SIDS. Quy mô doanh nghiệp dương nhưng không có ý nghĩa thống kê (β =
+0,082, p = 0,242).

\hypertarget{43-dux1ea1ng-uxe2m-ux111ux1a1n-ux111iux1ec7u-khuxf4ng-cuxf3-ux111ux1ed9-cong-h1-null-bux1eadc-hai}{%
\subsubsection{4.3 Dạng Âm Đơn Điệu: Không Có Độ Cong (H1, null bậc
hai)}\label{43-dux1ea1ng-uxe2m-ux111ux1a1n-ux111iux1ec7u-khuxf4ng-cuxf3-ux111ux1ed9-cong-h1-null-bux1eadc-hai}}

\textbf{M2} đưa vào hạng tử bậc hai FSTS\_c² để kiểm tra liệu quan hệ
I→P âm có cuối cùng bị đảo ngược hay không, tức là liệu có mô hình hình
chữ U ngược nào với phần dương nằm dưới phạm vi dữ liệu của chúng tôi
không.

Kết quả là: β(FSTS\_c) = −0,925, SE = 0,828, p = 0,265 (không có ý
nghĩa) β(FSTS\_c²) = +0,649, SE = 0,980, p = 0,508 (không có ý nghĩa)

Cả hạng tử tuyến tính lẫn bậc hai đều mất ý nghĩa thống kê khi được đưa
vào cùng nhau, cho thấy đa cộng tuyến cao giữa hai hạng tử trong mẫu
FSTS thấp này (FSTS trung bình = 0,060). Điều quan trọng là hệ số bậc
hai dương sẽ ngụ ý hình dạng chữ U thay vì hình chữ U ngược, một sự đảo
ngược của dạng thông thường mà bản thân nó nhất quán với FIP (quan hệ âm
đơn điệu, không có điểm ngưỡng đi lên trong phạm vi quan sát).

Theo Haans et al. (2016), chúng tôi đánh giá bốn điều kiện cần thiết cho
quan hệ có độ cong. Cả bốn điều kiện phải thỏa mãn để xác nhận hình chữ
U ngược; ở đây, cả bốn điều kiện đều \emph{thất bại}, đây là chữ ký kỳ
vọng của FIP:

\begin{itemize}
\tightlist
\item
  \textbf{C1} (độ dốc dương ở FSTS thấp): β(FSTS\_c) = −0,925 (\emph{p}
  = 0,265), \textbf{không có ý nghĩa thống kê và âm về dấu}; không có
  cánh tay đi lên ở các giá trị FSTS thấp
\item
  \textbf{C2} (độ cong âm β₂ \textless{} 0): β(FSTS\_c²) = +0,649
  (\emph{p} = 0,508), \textbf{không có ý nghĩa thống kê; dấu dương cho
  thấy hình chữ U, không phải hình chữ U ngược}
\item
  \textbf{C3} (điểm ngưỡng trong phạm vi dữ liệu): TP* = −(−0,925) / (2
  × 0,649) = 0,712, \textbf{điểm ngưỡng ngụ ý nằm ngoài mức tối đa mẫu
  quan sát (FSTS\_max ≈ 0,80, nhưng cả β₁ lẫn β₂ đều không có ý nghĩa
  thống kê)}
\item
  \textbf{C4} (kiểm định utest Lind--Mehlum, Lind \& Mehlum, 2010):
  utest \emph{p} = 0,941, \textbf{không bác bỏ giả thuyết null về quan
  hệ âm đơn điệu; không phát hiện hình chữ U ngược}
\end{itemize}

Mô hình là \textbf{quan hệ âm đơn điệu}, nhất quán với H1 và cơ chế cấu
trúc FIP. Không có điểm ngưỡng trong phạm vi FSTS quan sát được (0--1).
Điều này xác nhận FIP: không có điểm ngưỡng trong phạm vi dữ liệu.

\hypertarget{44-kiux1ec3m-ux111ux1ecbnh-vux1eefng}{%
\subsubsection{4.4 Kiểm Định
Vững}\label{44-kiux1ec3m-ux111ux1ecbnh-vux1eefng}}

Chúng tôi kiểm tra ba đặc tả bổ sung để đánh giá tính vững của kết quả
FIP.

\textbf{Đặc tả M\_yearFE (chỉ hiệu ứng cố định năm)}: Loại bỏ hiệu ứng
cố định quốc gia, có thể hấp thụ quá nhiều biến thiên cấu trúc, cho kết
quả tác động âm lớn hơn đáng kể:

β(FSTS) = \textbf{−1,236}, SE = 0,269, t = −4,594, p =
\textbf{\textless{} 0,001}

Hệ số lớn hơn này phản ánh biến thiên giữa các quốc gia bổ sung trong
quan hệ FSTS--hiệu quả. Kết quả FIP mạnh hơn đáng kể, xác nhận rằng bất
lợi cấu trúc cấp quốc gia làm trầm trọng thêm FIP cấp doanh nghiệp.

\textbf{Đặc tả M\_bivariate}: Hồi quy song biến của ln\_LP trên FSTS
(không có biến kiểm soát) cho kết quả:

β(FSTS) = \textbf{−1,596}, SE = 0,263, t = −6,062, p =
\textbf{\textless{} 0,001}

Mô hình âm đơn điệu vẫn tồn tại mà không có bất kỳ điều kiện hóa nào về
đặc điểm doanh nghiệp, tiếp tục xác nhận tính vững của FIP.

\textbf{Mẫu con chỉ nhà xuất khẩu} (N = 187): Trong số 187 nhà xuất khẩu
tích cực (FSTS \textgreater{} 0), tác động âm là:

β(FSTS\_c) = \textbf{−0,901}, SE = 0,398, p = \textbf{0,027}

Kết quả này rất quan trọng: ngay cả trong số các doanh nghiệp đã tham
gia xuất khẩu thành công, cường độ xuất khẩu cao hơn gắn liền với năng
suất thấp hơn. Điều này loại trừ giải thích lựa chọn đơn giản (ví dụ:
chỉ các doanh nghiệp năng suất thấp mới bị loại khỏi xuất khẩu, thúc đẩy
mối tương quan âm trong tổng hợp). Trong số các nhà xuất khẩu thực sự,
việc tăng cường xuất khẩu thêm làm xấu đi hiệu quả, đặc điểm của cơ chế
FIP chứ không phải lựa chọn bất lợi.

\hypertarget{45-kux1ebft-quux1ea3-null-nux103ng-lux1ef1c-h2}{%
\subsubsection{4.5 Kết Quả Null Năng Lực
(H2)}\label{45-kux1ebft-quux1ea3-null-nux103ng-lux1ef1c-h2}}

\textbf{M3} đưa vào TCI\_z và DAI\_z như là các biến hồi quy bổ sung để
kiểm định H2. Không có biến nào đạt được ý nghĩa thống kê:

β(TCI\_z) = +0,058, SE = 0,085, p = 0,495 (không có ý nghĩa) β(DAI\_z) =
+0,062, SE = 0,073, p = 0,402 (không có ý nghĩa)

Hệ số FIP (FSTS\_c) vẫn âm nhưng mất ý nghĩa thống kê trong M3 (β =
−0,974, p = 0,252), có khả năng phản ánh mẫu giảm đáng kể (N = 526 cho
M3, so với N = 1.469 cho M1, do các biến năng lực bị thiếu). Chiều hướng
của hệ số FSTS\_c vẫn âm và về số học lớn hơn trong M3, nhất quán với
FIP; sự mất ý nghĩa thống kê có thể quy cho việc giảm năng lực thống kê
từ hạn chế mẫu con.

Các kết quả này nhất quán với H2 (không điều tiết): Chỉ số năng lực công
nghệ (TCI\_z: β = +0,058, SE = 0,085, p = 0,495) và Chỉ số áp dụng số
(DAI\_z: β = +0,062, SE = 0,073, p = 0,402) không điều tiết có ý nghĩa
thống kê quan hệ I→P trong SIDS Thái Bình Dương. Điều này nhất quán với
cơ chế FIP áp đảo các tác động năng lực trong các nền kinh tế SIDS bị
ràng buộc về mặt cấu trúc.

Lưu ý quan trọng: Trong bối cảnh SIDS, DAI không phát huy hiệu quả do hạ
tầng số hạn chế. Chúng tôi không gọi đây là "năng lực số động" mà chỉ là
proxy sự hiện diện trang web nền tảng Tầng 1. Các ước lượng điểm dương
cho TCI\_z và DAI\_z (dù nhỏ và không có ý nghĩa thống kê) nhất quán với
diễn giải rằng các năng lực nâng cao mức năng suất chung. Diễn giải này
phù hợp với các bối cảnh châu Á đại lục: Goedhuys \& Sleuwaegen, 2013;
Teece, 2007) nhưng không thể điều chỉnh tính âm cấu trúc của quan hệ
FSTS--hiệu quả mà FIP mô tả.

\hypertarget{46-tuxf3m-tux1eaft-kux1ebft-quux1ea3}{%
\subsubsection{4.6 Tóm Tắt Kết
Quả}\label{46-tuxf3m-tux1eaft-kux1ebft-quux1ea3}}

Bảng 2 tóm tắt ma trận hệ số đầy đủ. Mô hình FIP được xác nhận qua bốn
điểm. (i) quan hệ tuyến tính FSTS--hiệu quả âm có ý nghĩa thống kê (M1:
β = −0,404, p = 0,032). (ii) không có độ cong bậc hai hay điểm ngưỡng có
ý nghĩa thống kê (M2: cả FSTS\_c lẫn FSTS\_c² đều không có ý nghĩa thống
kê). (iii) âm vững chắc trên bốn đặc tả (M\_yearFE: β = −1,236, p
\textless{} 0,001; M\_bivariate: β = −1,596, p \textless{} 0,001; chỉ
nhà xuất khẩu: β = −0,901, p = 0,027). (iv) điều tiết năng lực null (M3:
TCI\_z p = 0,495; DAI\_z p = 0,402).

\begin{center}\rule{0.5\linewidth}{0.5pt}\end{center}

\hypertarget{5-thux1ea3o-luux1eadn}{%
\subsection{5. Thảo Luận}\label{5-thux1ea3o-luux1eadn}}

\hypertarget{51-fip-nhux1b0-ux111iux1ec1u-kiux1ec7n-biuxean-cux1ee7a-muxf4-huxecnh-huxecnh-chux1eef-u-ngux1b0ux1ee3c}{%
\subsubsection{5.1 FIP như Điều Kiện Biên của Mô Hình Hình Chữ U
Ngược}\label{51-fip-nhux1b0-ux111iux1ec1u-kiux1ec7n-biuxean-cux1ee7a-muxf4-huxecnh-huxecnh-chux1eef-u-ngux1b0ux1ee3c}}

Phát hiện trung tâm của chúng tôi là quan hệ I→P hình chữ U ngược không
thỏa mãn trong SIDS Thái Bình Dương. Thay vì mô hình phi tuyến tính với
giai đoạn dương, chúng tôi quan sát quan hệ âm đơn điệu tồn tại qua bốn
đặc tả mô hình hóa và mạnh hơn trong số các nhà xuất khẩu thực sự so với
toàn mẫu. Điều này không nhất quán với dự báo hình chữ U ngược (Lu \&
Beamish, 2004) và đường cong S ba giai đoạn (Contractor et al., 2003),
cả hai đều dự báo ít nhất một phạm vi nào đó của quan hệ quốc tế
hóa--hiệu quả dương.

Khái niệm FIP giải thích sự đảo ngược này thông qua logic các điều kiện
tiền đề cấu trúc bị vi phạm. Khi không có thị trường nội địa khả thi,
FSTS không phải là biến lựa chọn chiến lược mà là sự cần thiết sinh tồn.
Các doanh nghiệp không thể tích lũy tài nguyên căn cứ nội địa và năng
lực tổ chức để tài trợ và duy trì quốc tế hóa có lợi nhuận. Khi chi phí
thương mại khuếch đại "chi phí kinh doanh" của doanh thu xuất khẩu lên
30--210\% (Winters \& Martins, 2004), kinh tế học biên của xuất khẩu là
âm ngay từ đầu. Khi khoảng trống thể chế là toàn diện (thiếu vắng các
trung gian tài trợ thương mại, thực thi hợp đồng không đáng tin cậy, hạ
tầng tiêu chuẩn hạn chế), các doanh nghiệp không thể chiếm đoạt lợi ích
học hỏi và quy mô mà lý thuyết quốc tế hóa hứa hẹn.

Điều này tái định khung hình chữ U ngược từ một quy luật thực nghiệm phổ
quát thành một quan hệ \emph{có điều kiện}: nó thỏa mãn khi và bởi vì
các điều kiện tiền đề cấu trúc được thỏa mãn. Phát hiện của văn liệu
hiện có rằng chất lượng thể chế \emph{điều tiết} quan hệ I→P (Marano et
al., 2016) là trường hợp đặc biệt của một quan hệ lý thuyết tổng quát
hơn. Khi các điều kiện thể chế vượt qua ngưỡng cấu trúc từ "điều tiết
yếu hơn" sang "điều kiện tiền đề bị vi phạm," dấu, không chỉ biên độ,
của quan hệ I→P đảo ngược.

\hypertarget{52-ux111ux1ecbnh-vux1ecb-luxfd-thuyux1ebft-tux1eeb-mux1ee9c-ux111ux1ed9-sang-loux1ea1i}{%
\subsubsection{5.2 Định Vị Lý Thuyết: Từ Mức Độ Sang
Loại}\label{52-ux111ux1ecbnh-vux1ecb-luxfd-thuyux1ebft-tux1eeb-mux1ee9c-ux111ux1ed9-sang-loux1ea1i}}

FIP đại diện cho sự chuyển dịch định tính trong cách chúng tôi khái niệm
hóa các tác động thể chế lên quan hệ I→P. Hầu hết các lý thuyết hóa
trước đây coi các tác động thể chế là các điều chỉnh liên tục của hình
dạng và vị trí đường cong hình chữ U ngược. Cụ thể, chất lượng thể chế
cao hơn dịch chuyển FSTS tối ưu lên phía trên, mở rộng phạm vi dương,
hoặc nâng cao mức hiệu quả đỉnh (Marano et al., 2016; Hitt et al.,
1997). Đây là sự thay đổi về \emph{mức độ}.

Các phát hiện của chúng tôi chỉ ra một sự thay đổi về \emph{loại}: ở cực
đoan của phổ thể chế, nơi tất cả ba điều kiện tiền đề cấu trúc đồng thời
vắng mặt, hình thức hàm I→P không phải là hình chữ U ngược được sửa đổi
mà là một mô hình định tính khác (quan hệ âm đơn điệu). Điều này mở rộng
vốn từ lý thuyết của nghiên cứu IB vượt ra ngoài "điều tiết cường độ"
sang "điều tiết sự tồn tại" của quan hệ I→P.

Phát hiện này kết nối với câu hỏi rộng hơn về điều kiện biên trong lý
thuyết IB (Verbeke \& Brugman, 2018; Hennart, 2007). Nếu các lý thuyết
I→P chính lưu được rút ra từ các mẫu nền kinh tế nơi các điều kiện tiền
đề cấu trúc được thỏa mãn rộng rãi, những lý thuyết đó có thể không áp
dụng một cách hệ thống cho bối cảnh SIDS, không chỉ cần điều chỉnh, mà
đại diện cho một chế độ lý thuyết định tính khác. FIP cung cấp một nhãn
hiệu có cơ sở lý thuyết cho chế độ thay thế đó.

\hypertarget{53-null-nux103ng-lux1ef1c-quyux1ebft-ux111ux1ecbnh-thux1ec3-chux1ebf-vs-ux111iux1ec1u-tiux1ebft}{%
\subsubsection{5.3 Null Năng Lực: Quyết Định Thể Chế vs. Điều
Tiết}\label{53-null-nux103ng-lux1ef1c-quyux1ebft-ux111ux1ecbnh-thux1ec3-chux1ebf-vs-ux111iux1ec1u-tiux1ebft}}

Kết quả null cho TCI\_z và DAI\_z trong M3 thêm chiều thứ hai vào lý
thuyết FIP. Trong các bối cảnh châu Á đại lục, Việt Nam, Singapore,
Trung Quốc, và các mẫu đa quốc gia, Chỉ số năng lực công nghệ (TCI) và
Chỉ số áp dụng số (DAI) là các biến điều tiết có ý nghĩa thống kê của
quan hệ I→P (Goedhuys \& Sleuwaegen, 2013; Teece, 2007). Tại sao chúng
null trong SIDS?

Cơ chế FIP cung cấp câu trả lời: các năng lực điều tiết \emph{chi phí
phức tạp tổ chức} ở phần giảm dần của hình chữ U ngược, nhưng cơ chế FIP
hoạt động ở mức độ cấu trúc, tiền-tổ chức. Một doanh nghiệp SIDS với
chứng nhận ISO 9001 (TCI cao) vẫn phải đối mặt với mức khuếch đại chi
phí thương mại 30--210\% như nhau; một doanh nghiệp SIDS có trang web
hoạt động (DAI = 1) vẫn phải đối mặt với các thể chế tài trợ thương mại
vắng mặt. Các năng lực là các biến điều tiết mạnh trong khuôn khổ điều
kiện tiền đề cấu trúc trở nên kém hiệu quả khi những điều kiện tiền đề
đó không thỏa mãn.

Điều này nhất quán với sự nhấn mạnh của quan điểm dựa trên nguồn lực về
tính phụ thuộc bối cảnh (Barney, 1991): các nguồn lực tạo ra giá trị chỉ
trong các môi trường thể chế hỗ trợ khả năng chiếm đoạt. Khung năng lực
động của Teece (2007) tương tự giả định các cơ chế thị trường hoạt động;
khi vắng mặt chúng, các năng lực động có thể bị "mắc kẹt" trong doanh
nghiệp, không thể chuyển hóa thành lợi thế hiệu quả thông qua các giao
dịch thị trường quốc tế.

Các ước lượng điểm dương, mặc dù không có ý nghĩa thống kê, cho TCI\_z
và DAI\_z gợi ý rằng các năng lực nâng cao mức năng suất chung (nhất
quán với vai trò của chúng trong văn liệu rộng hơn), nhưng không thể
điều chỉnh tính âm hướng của quan hệ FSTS--hiệu quả trong bối cảnh cấu
trúc SIDS. Sự tinh tế này, các năng lực giúp mức, không phải độ dốc, có
thể là một hướng nghiên cứu SIDS trong tương lai có tiềm năng.

\hypertarget{54-huxe0m-uxfd-chuxednh-suxe1ch-chux1b0ux1a1ng-truxecnh-nghux1ecb-sux1ef1-antigua-vuxe0-barbuda}{%
\subsubsection{5.4 Hàm Ý Chính Sách: Chương Trình Nghị Sự Antigua và
Barbuda}\label{54-huxe0m-uxfd-chuxednh-suxe1ch-chux1b0ux1a1ng-truxecnh-nghux1ecb-sux1ef1-antigua-vuxe0-barbuda}}

Các phát hiện của chúng tôi có hàm ý trực tiếp cho chính sách phát triển
SIDS. Chương trình nghị sự Antigua và Barbuda dành cho SIDS (United
Nations General Assembly, 2024), được thông qua tại Hội nghị Quốc tế thứ
tư về SIDS, đặt thúc đẩy xuất khẩu và thuận lợi hóa thương mại là các
trụ cột trung tâm của chiến lược phát triển SIDS. Cách đặt vấn đề này
ngầm định giả định rằng tăng FSTS sẽ cải thiện hiệu quả doanh nghiệp và
qua đó kích thích tăng trưởng kinh tế.

Kết quả của chúng tôi cảnh báo chống lại giả định này. Nếu FIP là đặc
điểm cấu trúc của nền kinh tế doanh nghiệp SIDS Thái Bình Dương, thì các
chính sách đẩy doanh nghiệp SIDS hướng đến cường độ xuất khẩu cao hơn
(mà không đồng thời giải quyết các điều kiện tiền đề cấu trúc bị vi
phạm) có thể làm xấu đi hiệu quả doanh nghiệp ở cấp vi mô ngay cả khi
cải thiện cán cân thương mại kinh tế vĩ mô. Câu hỏi chính sách liên quan
trở thành: những can thiệp nào có thể giải quyết ba điều kiện tiền đề
cấu trúc tạo ra FIP?

Về \textbf{Điều kiện tiền đề 1} (thị trường nội địa khả thi): các sáng
kiến hội nhập khu vực (Hiệp định Kinh tế Quan hệ Gần gũi Hơn Thái Bình
Dương Plus, PACER Plus) tạo ra quyền tiếp cận thị trường hiệu quả với
các SIDS lân cận và với Australia và New Zealand. Cách tiếp cận này có
thể mở rộng thị trường nội địa hiệu quả vượt ra ngoài biên giới quốc
gia, giảm bắt buộc cấu trúc phải xuất khẩu sang các thị trường xa xôi
chi phí cao.

Về \textbf{Điều kiện tiền đề 2} (chi phí thương mại có thể quản lý
được): giảm chi phí vận tải thông qua các chương trình kết nối hàng hải
khu vực và, nơi khả thi, trợ cấp hàng không cho xuất khẩu giá trị
cao-trọng lượng thấp (nông sản, thủ công mỹ nghệ, dịch vụ) có thể giảm
mức khuếch đại chi phí thương mại thúc đẩy FIP.

Về \textbf{Điều kiện tiền đề 3} (thể chế có chức năng): phát triển thể
chế có mục tiêu tập trung vào hạ tầng thuận lợi hóa thương mại, các cơ
quan tiêu chuẩn, trung gian tài trợ thương mại, cơ quan xúc tiến xuất
khẩu với năng lực thực sự, thay vì cải cách quản trị chung có thể giải
quyết trực tiếp nhất cơ chế FIP. Bằng chứng của Goedhuys và Sleuwaegen
(2013) về chứng nhận quốc tế rất có tính hướng dẫn: phát hiện của họ
rằng chứng nhận ISO cải thiện năng suất doanh nghiệp nhất quán với kết
quả TCI null ở đây nếu lợi ích năng suất của chứng nhận có điều kiện
trên các thể chế hỗ trợ thực thi tiêu chuẩn và nhận biết khách hàng, các
điều kiện vắng mặt trong bối cảnh SIDS.

\hypertarget{55-hux1ea1n-chux1ebf}{%
\subsubsection{5.5 Hạn Chế}\label{55-hux1ea1n-chux1ebf}}

Một số hạn chế giới hạn diễn giải các phát hiện của chúng tôi.

Thứ nhất, \textbf{cấu trúc tổng hợp chéo tiết diện}: mặc dù chúng tôi
tổng hợp trên nhiều đợt quốc gia-năm (2009--2025), ít nền kinh tế SIDS
có nhiều hơn một đợt WBES (xem Bảng 1), giới hạn nhận dạng bảng dữ liệu
của các tác động nhân quả. Ước lượng FIP chủ yếu được xác định từ biến
thiên chéo tiết diện trong FSTS trong các ô quốc gia-năm. Nghiên cứu
tương lai sử dụng dữ liệu bảng, có sẵn cho Fiji và Papua New Guinea với
hai đợt, có thể tăng cường nhận dạng nhân quả.

Thứ hai, \textbf{đo lường FSTS trong bối cảnh xuất khẩu thưa thớt}: với
FSTS trung bình 0,060 và chỉ 12,7\% nhà xuất khẩu, phân phối FSTS có
đuôi phải rất dài. Các ước lượng hệ số trong M2 có thể nhạy cảm với một
số giá trị ngoại lệ FSTS cao; kết quả bậc hai không có ý nghĩa thống kê
cần được diễn giải cẩn thận. Nghiên cứu tương lai có thể sử dụng hồi quy
phân vị hoặc winsorization để đánh giá độ nhạy phân phối.

Thứ ba, \textbf{DAI chỉ là proxy Tầng 1}: thước đo áp dụng số của chúng
tôi chỉ nắm bắt sự hiện diện trang web, một chỉ số nền tảng, tĩnh. Điều
này phù hợp với bối cảnh SIDS do tính sẵn có hạn chế của các chỉ số hạ
tầng số tiên tiến hơn trong WBES. Tuy nhiên, nó có nghĩa là chúng tôi
không thể đánh giá liệu các năng lực số tinh vi hơn (thương mại điện tử,
thanh toán di động, điện toán đám mây) có thể điều tiết FIP theo cách
khác hay không. Chúng tôi rõ ràng không nhận định rằng mình đang đo năng
lực số động.

Thứ tư, \textbf{khả năng tổng quát hóa ngoài SIDS Thái Bình Dương}: FIP
được lý thuyết hóa là phát sinh khi ba điều kiện tiền đề cấu trúc đồng
thời vắng mặt. Các bối cảnh cực đoan khác, các vi quốc gia Sahara, các
nền kinh tế cao nguyên xa xôi, các bối cảnh hậu xung đột, có thể thể
hiện các mô hình tương tự. SIDS Caribbean (không được đưa vào mẫu của
chúng tôi) đại diện cho bối cảnh mở rộng tự nhiên.

\begin{center}\rule{0.5\linewidth}{0.5pt}\end{center}

\hypertarget{6-kux1ebft-luux1eadn}{%
\subsection{6. Kết Luận}\label{6-kux1ebft-luux1eadn}}

Chúng tôi giới thiệu Gánh nặng quốc tế hóa bắt buộc (Forced
Internationalization Penalty, FIP). Đây là tác động âm đơn điệu của
cường độ quốc tế hóa lên hiệu quả doanh nghiệp khi ba điều kiện tiền đề
cấu trúc của quan hệ hình chữ U ngược đồng thời vắng mặt: thị trường nội
địa khả thi, chi phí thương mại có thể quản lý được, và hỗ trợ thể chế
có chức năng.

Nghiên cứu sử dụng dữ liệu WBES từ 1.469 doanh nghiệp thuộc chín nền
kinh tế SIDS Thái Bình Dương. Kết quả tìm thấy hỗ trợ vững chắc cho FIP
trên bốn đặc tả thực nghiệm (β = −0,404, p = 0,032 trong mô hình hiệu
ứng cố định quốc gia-năm chính; β = −1,236, p \textless{} 0,001 không có
hiệu ứng cố định quốc gia; β = −0,901, p = 0,027 trong mẫu con chỉ nhà
xuất khẩu). Hạng tử bậc hai không có ý nghĩa thống kê, xác nhận dạng âm
đơn điệu, không có điểm ngưỡng trong phạm vi quan sát. Chỉ số năng lực
công nghệ (TCI) và Chỉ số áp dụng số (DAI), vốn có ý nghĩa thống kê
trong các bối cảnh châu Á đại lục, lại null trong môi trường SIDS, nhất
quán với cơ chế lý thuyết FIP.

Các phát hiện này đóng góp vào lý thuyết IB theo hai cách. Thứ nhất,
chúng giúp định hướng lại diễn ngôn về các tác động thể chế từ "điều
tiết thể chế cường độ I→P" sang "quyết định thể chế dấu I→P." Hình chữ U
ngược có thể có điều kiện trên các điều kiện tiền đề cấu trúc thay vì áp
dụng rộng rãi trên các bối cảnh. Thứ hai, chúng mở rộng ứng dụng "quốc
tế hóa bắt buộc" vượt ra ngoài văn liệu SME sang lĩnh vực kết quả hiệu
quả. Đóng góp này cung cấp một cấu trúc có định nghĩa lý thuyết và cơ sở
thực nghiệm (FIP) mà nghiên cứu trong tương lai có thể xây dựng trên.

Đối với chương trình nghị sự chính sách phát triển thế giới, FIP gợi ý
rằng các chiến lược thúc đẩy xuất khẩu trong SIDS đòi hỏi các can thiệp
cấu trúc bổ sung. Những can thiệp này cần giải quyết chi phí thương mại,
khả năng khả thi của thị trường nội địa thông qua hội nhập khu vực, và
phát triển thể chế có mục tiêu, chứ không chỉ tăng cường hoạt động xuất
khẩu thuần túy.

\begin{center}\rule{0.5\linewidth}{0.5pt}\end{center}

\hypertarget{tuxe0i-liux1ec7u-tham-khux1ea3o}{%
\subsection{Tài Liệu Tham Khảo}\label{tuxe0i-liux1ec7u-tham-khux1ea3o}}

Barney, J. (1991). Firm resources and sustained competitive advantage.
\emph{Journal of Management, 17}(1), 99--120.

Bausch, A., \& Krist, M. (2007). The effect of context-related
moderators on the internationalization-performance relationship.
\emph{Management International Review, 47}(3), 319--347.

Briguglio, L. (1995). Small island developing states and their economic
vulnerabilities. \emph{World Development, 23}(9), 1615--1632.

Briguglio, L., Cordina, G., Farrugia, N., \& Vella, S. (2009). Economic
vulnerability and resilience: Concepts and measurements. \emph{Oxford
Development Studies, 37}(3), 229--247.

Contractor, F. J., Kundu, S. K., \& Hsu, C.-C. (2003). A three-stage
theory of international expansion: The link between multinationality and
performance in the service sector. \emph{Journal of International
Business Studies, 34}(1), 5--18.

Diamantopoulos, A., \& Winklhofer, H. M. (2001). Index construction with
formative indicators: An alternative to scale development. \emph{Journal
of Marketing Research, 38}(2), 269--277.

Doh, J. P., Rodrigues, S., Saka-Helmhout, A., \& Makhija, M. (2017).
International business responses to institutional voids. \emph{Journal
of International Business Studies, 48}(3), 293--307.

Goedhuys, M., \& Sleuwaegen, L. (2013). The impact of international
standards certification on the performance of firms in less developed
countries. \emph{World Development, 46}, 87--101.

Haans, R. F. J., Pieters, C., \& He, Z.-L. (2016). Thinking about U:
Theorizing and testing U- and inverted U-shaped relationships in
strategy research. \emph{Strategic Management Journal, 37}(7),
1177--1195. \url{https://doi.org/10.1002/smj.2399}

Hennart, J.-F. (2007). The theoretical rationale for a
multinationality-performance relationship. \emph{Management
International Review, 47}(3), 423--452.

Hitt, M. A., Hoskisson, R. E., \& Kim, H. (1997). International
diversification: Effects on innovation and firm performance in
product-diversified firms. \emph{Academy of Management Journal, 40}(4),
767--798.

Jarvis, C. B., MacKenzie, S. B., \& Podsakoff, P. M. (2003). A critical
review of construct indicators and measurement model misspecification in
marketing and consumer research. \emph{Journal of Consumer Research,
30}(2), 199--218.

Johanson, J., \& Vahlne, J.-E. (1977). The internationalization process
of the firm. \emph{Journal of International Business Studies, 8}(1),
23--32.

Johanson, J., \& Vahlne, J.-E. (2009). The Uppsala internationalization
process model revisited. \emph{Journal of International Business
Studies, 40}(9), 1411--1431.

Khanna, T., \& Palepu, K. G. (2010). \emph{Winning in emerging markets:
A road map for strategy and execution}. Harvard Business Press.

Kirca, A. H., Hult, G. T. M., Roth, K., Cavusgil, S. T., Perry, M. Z.,
Akdeniz, M. B., Deligonul, S. Z., Mena, J. A., Pollitte, W. A., Hult, R.
N., White, J. C., \& White, R. C. (2011). Firm-specific assets,
multinationality, and financial performance: A meta-analytic review and
theoretical integration. \emph{Academy of Management Journal, 54}(1),
47--72.

Lall, S. (1992). Technological capabilities and industrialization.
\emph{World Development, 20}(2), 165--186.

Lind, J. T., \& Mehlum, H. (2010). With or without U? The appropriate
test for a U-shaped relationship. \emph{Oxford Bulletin of Economics and
Statistics, 72}(1), 109--118.

Lu, J. W., \& Beamish, P. W. (2004). International diversification and
firm performance: The S-curve hypothesis. \emph{Academy of Management
Journal, 47}(4), 598--609.

Mahler, D., Serajuddin, U., Wadhwa, D., \& Yonzan, N. (2026, May).
Global development: Progress and challenges. In \emph{Atlas of Global
Development 2026}. World Bank Data360. World Bank Group.

Mahler, D., Wang, H., \& Weber, M. (2026, May). Inequalities in use of
and exposure to artificial intelligence. In \emph{Atlas of Global
Development 2026}. World Bank Data360. World Bank Group.
\url{https://data360.worldbank.org/en/int/atlas/artificial-intelligence}

Marano, V., Tallman, S., \& Teece, D. J. (2016). The context of
innovation and the innovation of context. \emph{Journal of Leadership \&
Organizational Studies, 23}(4), 421--428.

North, D. C. (1990). \emph{Institutions, institutional change and
economic performance}. Cambridge University Press.

Pirlea, A. F., Wadhwa, D., Mahler, D., Serajuddin, U., Welch, M., Thudt,
A., \& Lambrechts, M. (2026, May). Global progress. In \emph{Atlas of
Global Development 2026}. World Bank Data360. World Bank Group.
\url{https://data360.worldbank.org/en/int/atlas/global-progress}

Read, R. (2004). The implications of increasing globalization and
regionalism for the economic growth of small island states. \emph{World
Development, 32}(2), 365--378.

Sarasvathy, S. D., Kumar, K., York, J. G., \& Bhagavatula, S. (2014). An
effectual approach to international entrepreneurship: Overlaps,
challenges, and provocative possibilities. \emph{Entrepreneurship Theory
and Practice, 38}(1), 71--93.

Scott, W. R. (1995). \emph{Institutions and organizations}. Sage.

Teece, D. J. (2007). Explicating dynamic capabilities: The nature and
microfoundations of (sustainable) enterprise performance.
\emph{Strategic Management Journal, 28}(13), 1319--1350.

United Nations General Assembly. (2024). \emph{The Antigua and Barbuda
Agenda for Small Island Developing States: A renewed declaration for
resilient prosperity} (A/RES/78/317). United Nations.

Vahlne, J.-E., \& Johanson, J. (2020). The Uppsala model: Networks and
micro-foundations. \emph{Journal of International Business Studies,
51}(1), 4--10.

Verbeke, A., \& Brugman, M. (2018). Triple-testing the quality of
multinationality--performance research: An internalization theory
perspective. \emph{Journal of International Business Studies, 49}(7),
801--822.

Winters, L. A., \& Martins, P. M. G. (2004). When comparative advantage
is not enough: Business costs in small remote economies. \emph{World
Trade Review, 3}(3), 347--383.

Zaheer, S. (1995). Overcoming the liability of foreignness.
\emph{Academy of Management Journal, 38}(2), 341--363.

\begin{center}\rule{0.5\linewidth}{0.5pt}\end{center}

\hypertarget{phux1ee5-lux1ee5c-a-bux1ea3ng-sux1ed1-liux1ec7u}{%
\subsection{Phụ Lục A: Bảng Số
Liệu}\label{phux1ee5-lux1ee5c-a-bux1ea3ng-sux1ed1-liux1ec7u}}

\hypertarget{bux1ea3ng-1-thux1ed1ng-kuxea-tuxf3m-tux1eaft-cux1ea5p-quux1ed1c-gia}{%
\subsubsection{Bảng 1. Thống Kê Tóm Tắt Cấp Quốc
Gia}\label{bux1ea3ng-1-thux1ed1ng-kuxea-tuxf3m-tux1eaft-cux1ea5p-quux1ed1c-gia}}

\begin{longtable}[]{@{}llllll@{}}
\toprule\noalign{}
Nền kinh tế & N (doanh nghiệp) & Nhà xuất khẩu & Tỷ lệ xuất khẩu (\%) &
FSTS trung bình & Đợt khảo sát \\
\midrule\noalign{}
\endhead
\bottomrule\noalign{}
\endlastfoot
Fiji & --- & --- & --- & --- & --- \\
Kiribati & --- & --- & --- & --- & --- \\
Papua New Guinea & --- & --- & --- & --- & --- \\
Samoa & --- & --- & --- & --- & --- \\
Solomon Islands & --- & --- & --- & --- & --- \\
Timor-Leste & --- & --- & --- & --- & --- \\
Tonga & --- & --- & --- & --- & --- \\
Vanuatu & --- & --- & --- & --- & --- \\
Comoros & --- & --- & --- & --- & --- \\
\textbf{Tổng cộng} & \textbf{1.469} & \textbf{187} & \textbf{12,7\%} &
\textbf{0,060} & 2009--2025 \\
\end{longtable}

\emph{Ghi chú: Số liệu phân tách cấp quốc gia có sẵn trong kho lưu trữ
nhân rộng (\texttt{p8/replication/}). Tổng N phản ánh mẫu phân tích
(không thiếu ln\_LP và FSTS). Phạm vi đợt khảo sát khác nhau theo quốc
gia.}

\hypertarget{bux1ea3ng-2-kux1ebft-quux1ea3-hux1ed3i-quy-cux1b0ux1eddng-ux111ux1ed9-quux1ed1c-tux1ebf-huxf3a-vuxe0-nux103ng-suux1ea5t-lao-ux111ux1ed9ng}{%
\subsubsection{Bảng 2. Kết Quả Hồi Quy: Cường Độ Quốc Tế Hóa và Năng
Suất Lao
Động}\label{bux1ea3ng-2-kux1ebft-quux1ea3-hux1ed3i-quy-cux1b0ux1eddng-ux111ux1ed9-quux1ed1c-tux1ebf-huxf3a-vuxe0-nux103ng-suux1ea5t-lao-ux111ux1ed9ng}}

\begin{longtable}[]{@{}lllllll@{}}
\toprule\noalign{}
& M0 Kiểm soát & M1 Tuyến tính & M2 Bậc hai & M3 Năng lực & M\_yearFE &
M\_bivariate \\
\midrule\noalign{}
\endhead
\bottomrule\noalign{}
\endlastfoot
\textbf{FSTS\_c} & & −0,404* (0,188) & −0,925 (0,828) & −0,974 (0,850) &
& \\
\textbf{FSTS\_c²} & & & +0,649 (0,980) & +0,782 (1,001) & & \\
\textbf{FSTS} & & & & & −1,236*** (0,269) & −1,596*** (0,263) \\
\textbf{TCI\_z} & & & & +0,058 (0,085) & & \\
\textbf{DAI\_z} & & & & +0,062 (0,073) & & \\
ln\_size & +0,080 (0,070) & +0,082 (0,070) & +0,084 (0,070) & +0,062
(0,070) & +0,247* (0,097) & --- \\
firm\_age & +0,009* (0,005) & +0,009. (0,005) & +0,009. (0,005) &
+0,010* (0,005) & +0,047*** (0,009) & --- \\
foreign\_own\_pct & −0,006. (0,003) & −0,006. (0,003) & −0,006. (0,003)
& −0,006. (0,003) & −0,006* (0,003) & --- \\
Hiệu ứng cố định Quốc gia & ✓ & ✓ & ✓ & ✓ & --- & --- \\
Hiệu ứng cố định Năm & ✓ & ✓ & ✓ & ✓ & ✓ & --- \\
N & 1.469 & 1.469 & 1.469 & 526 & 1.469 & 1.469 \\
R² & 0,799 & 0,800 & 0,800 & --- & 0,512 & --- \\
\end{longtable}

\emph{Ghi chú: Hệ số được báo cáo với sai số chuẩn vững với phương sai
không đồng nhất (HC1) trong ngoặc đơn. Hiệu ứng cố định quốc gia được
hấp thụ trong M0--M3. *** p \textless{} 0,001; ** p \textless{} 0,01; *
p \textless{} 0,05; . p \textless{} 0,10.}

\emph{Gánh nặng quốc tế hóa bắt buộc (FIP): Đã xác nhận, β(FSTS\_c) =
−0,404, p = 0,032 (M1).}

\hypertarget{bux1ea3ng-3-kiux1ec3m-ux111ux1ecbnh-vux1eefng-mux1eabu-con-chux1ec9-nhuxe0-xuux1ea5t-khux1ea9u-n--187}{%
\subsubsection{Bảng 3. Kiểm Định Vững: Mẫu Con Chỉ Nhà Xuất Khẩu (N =
187)}\label{bux1ea3ng-3-kiux1ec3m-ux111ux1ecbnh-vux1eefng-mux1eabu-con-chux1ec9-nhuxe0-xuux1ea5t-khux1ea9u-n--187}}

\begin{longtable}[]{@{}lllll@{}}
\toprule\noalign{}
Đặc tả & β(FSTS\_c) & SE & Giá trị p & Chiều hướng \\
\midrule\noalign{}
\endhead
\bottomrule\noalign{}
\endlastfoot
OLS, hiệu ứng cố định quốc gia+năm, N=187 & −0,901 & 0,398 & 0,027* &
Âm \\
M1 toàn mẫu & −0,404 & 0,188 & 0,032* & Âm \\
M\_yearFE toàn mẫu & −1,236 & 0,269 & \textless0,001*** & Âm \\
M\_bivariate & −1,596 & 0,263 & \textless0,001*** & Âm \\
\end{longtable}

\emph{Ghi chú: Gánh nặng được xác nhận trên tất cả các đặc tả. Kết quả
chỉ nhà xuất khẩu (β = −0,901) giải quyết lo ngại lựa chọn bất lợi:
trong số các nhà xuất khẩu tích cực, tăng cường xuất khẩu thêm làm xấu
đi hiệu quả.}

\begin{center}\rule{0.5\linewidth}{0.5pt}\end{center}

\emph{Hết bản dịch tiếng Việt}

\emph{Tương ứng với: p8/p8\_pacific\_sids\_en\_clean.md (394 dòng)}
\emph{Bản dịch này tuân theo chuẩn thuật ngữ CTU
(09b\_vn\_term\_glossary.md v1.4)} \emph{Phân loại ICRV: Nhóm VI ICRV
(SIDS/Frontier), nhóm thể chế yếu nhất/nhỏ nhất}

\end{document}
